\begin{proof}
Denote $\Delta=\cl(G^c)$ and $\Delta'=\cl(G'^c)$.
Let $A=\Vr(\bar{\St}_a (\Delta)),B=\Vr(\bar{\St}_b(\Delta)), C=A\cap B$ and $D=\Vr(\Delta)-(A\cup B)$. With this notation we observe the following crucial relation between $\Delta$ and $\Delta'$:
\begin{equation}\label{eq:Delta'}
 \Delta'=(\Delta[A]\cup_{\Delta[C]} \Delta[B])\cup_{\Delta[C]} (\Delta[C] \cup_{d\in D} (\{d\}*\Delta[C \cap \Vr(\St_d (\Delta))])).
\end{equation}

 We obtain this equality simply from the definition of $G'$, where every neighbour of $a$ is connected to every neighbour of $b$.
 The above decomposition of $\Delta'$ shows us how to prove the theorem using Hochster's formula and Mayer--Vietoris sequences, detailed next.


 Let $\reg(I(G')) = l+2$ and $W$ be a subset of the vertices of minimal size such that $\widetilde{H}_l (\Delta'[W]) \neq 0$.
 Decompose $W=W_1 \cup_{W_C} W_2 \subseteq V$ where $W_1=W\cap(A\cup_C B)$ (recall $C=A\cap B$), $W_2=W\cap(C\cup D)$ and $W_C=W_1\cap W_2=C \cap W$.

If $W\cap D=\emptyset$ then Eq.(\ref{eq:Delta'}) implies $\Delta'[W]=\Delta[W]$ and thus $\reg(I(G))\ge l+2$ as claimed.
Thus, we may assume
$W\cap D\neq \emptyset$, so $|W_1|<|W|$ and thus, by minimality of $W$, $\widetilde{H}_l(\Delta'[W_1])=0$.

If $|W_2|<|W|$ then minimality of $|W|$ also imply $\widetilde{H}_l(\Delta'[W_2])=0$.
By~\eqref{eq:Delta'} the decomposition $\Delta'[W]=\Delta'[W_1] \cup_{\Delta'[W_C]} \Delta'[W_2]$ holds. Consider the following Mayer--Vietoris exact sequence: $$\widetilde{H}_l \Delta'[W_1] \bigoplus \widetilde{H}_l \Delta'[W_2] \longrightarrow \widetilde{H}_l \Delta'[W] \longrightarrow \widetilde{H}_{l-1} (\Delta'[W_C]).$$
As the two summands on the left term vanish, exactness implies
$\widetilde{H}_{l-1} (\Delta'[W_C])\neq 0$. 
Now by~\eqref{eq:Delta'} we have  $\Delta'[W_C]=\Delta[W_C]$, hence suspension gives
  $0 \neq \widetilde{H}_l(\Delta[W_C] * \{\{a\},\{b\},\emptyset\})= \widetilde{H}_l(\Delta[\{a,b\} \cup W_C])$.
 We conclude $\reg (\Delta) \geq l+2$ as claimed.

Thus we may assume additionally that
$W_2=W$. 
Denote $X=\Delta'[W]=\Delta'[W_2]$ and let $d\in W\cap D$.
By~\eqref{eq:Delta'} $\text{link}_d X$ is an induced subcomplex of $\Delta[C]$, and further, the suspension of $\text{link}_d X$ by the vertices $a$ and $b$ is an induced
subcomplex of $\Delta$.

Thus, if
$\widetilde{H}_{l-1}(\text{link}_d X)\neq 0$
then
$0\neq \widetilde{H}_l (\{\{a\},\{b\},\emptyset\}*\text{link}_d X)=\widetilde{H}_l(\Delta[\{a,b\}\cup V(\text{link}_d X)])$, and thus $\reg(I(G))\ge l+2$ as claimed. 
So we may assume $\widetilde{H}_{l-1}(\text{link}_d X)=0$.


Consider the Mayer--Vietoris long exact sequence
corresponding to the union
$X=\text{ast}_d X\cup_{\text{link}_d X} \bar{\text{St}}_d X$:
$$\widetilde{H}_l (\text{ast}_{d} X) \bigoplus \widetilde{H}_l (\overline{\text{St}}_{d} X) \longrightarrow \widetilde{H}_l (X) \longrightarrow \widetilde{H}_{l-1} (\text{link}_d X).$$
By assumption, the right term is zero while the middle term is nonzero, and closed stars have vanishing homology, hence
$\widetilde{H}_l (\text{ast}_{d} X)=
\widetilde{H}_l (\Delta'[W\setminus\{d\}])\neq 0$.
This contradicts the minimality of $W$, so in fact this last case cannot occur.
\end{proof}
